 %!TEX root = main.tex


% ── Abstract ──────────────────────────────────────────────────────────────────
\twocolumn[
  \vspace{1em}
  {\small
  \begin{center}
    \textbf{\abstractname}
  \end{center}
  \begin{quotation}
This report presents the design and stress-testing of the GlobalFlow international
supply chain network, which distributes three product families — fertilisers~(A),
semiconductors~(B), and battery components~(C) — across a five-hub, fifteen-warehouse
multi-echelon structure connecting nine suppliers to thirty customers.
Phase~1 solves a Mixed-Integer Program (MIP) to identify the cost-minimising baseline
configuration, yielding an optimal total cost of $Z^* = 4,454,800$, with
14 warehouses open and 11 optional arcs activated.
Phase~2 subjects this network to six adversarial scenarios combining tariff shocks,
energy crises, hub closures, and corridor blockades.
For each scenario, three response strategies are evaluated: pure rerouting~(R),
network adaptation~(A), and full redesign~(F).
A central finding is that Scenarios S1 and S3 — involving the closure of
hub~H3 (Singapore) — render the standard formulation structurally infeasible,
because suppliers S4, S5, S6, and S9 are exclusively connected through H3,
stranding 100\% of semiconductor supply and 33\% of battery component supply.
We address this with a modified formulation that introduces unmet-demand slack
variables for strategies~R and~A, and emergency sourcing arcs for strategy~F.
Strategic analysis reveals a persistent single-point-of-failure in the baseline
topology and quantifies the cost of building a more resilient network.
  \end{quotation}}
  \vspace{1.5em}
]
